\chapter{The question, the biology, and the symbols}
\label{chap:foundations}

\section{What are we trying to understand?}

Imagine a small population. Most individuals carry one genetic type. A few
individuals carry another type. The second type gives its carriers a
reproductive advantage. Does that advantage guarantee that the second type
will spread through the whole population?

No. An advantage changes the chances of reproduction. It does not determine
the outcome of each reproduction event. A rare type can disappear before its
advantage has much opportunity to act.

Our question is therefore precise: \emph{what is the probability that an
initially rare, advantageous type eventually occupies the whole population?}
We will answer this question in two ways. First, we will repeat a random
computer experiment. Second, we will calculate an exact mathematical answer
for the model. Comparing the two answers will help us check our code.

The model is deliberately small. Its purpose is to make the reasoning visible.
It is not a claim that every biological population follows these rules.

\section{The biological words, one at a time}

\subsection{Population, individual, and genetic type}

A \textbf{population} is the collection of individuals included in our model.
An \textbf{individual} is one member of that collection. We follow one genetic
difference. An \textbf{allele} is a version of a genetic locus, meaning a
particular location or inherited genetic feature under consideration.

We call the two versions type A and type B. These letters are labels. We do
not need to name a real gene. Type A is the \emph{focal} type: it is the type
whose fate we follow. Type B is the \emph{resident} type: it is the alternative
already present in the population.

We use a \textbf{haploid} model. Each individual carries one copy of the locus
we follow, so each individual is either A or B. We do not have an AB genotype
in this model. We do not model diploid inheritance, dominance, or mating.
This is a modelling choice, not a description of all organisms.

Sometimes we call A the \emph{mutant} type. Here that word describes the type
whose initial arrival we imagine. It does not mean new mutations continue to
appear during a run. After setting the initial state, we allow no further
mutation.

\subsection{Loss and fixation}

Type A is \textbf{lost} when no A individuals remain. It has reached
\textbf{fixation} when every individual is A.

These outcomes concern this modelled population and this locus. Fixation
does not mean that all organisms everywhere carry A. Loss of A does not mean
that the whole population has died. The population still exists; it now
contains only B individuals.

Because we allow no new mutation, A cannot return after it has been lost.
Likewise, B cannot return after A has fixed. These are \textbf{absorbing
states}: once the process enters one, it stays there.

\subsection{Selection and drift}

\textbf{Selection} means that the genetic types differ in their expected
reproductive success. In this project, A can have a greater chance of being
chosen as a parent than B.

\textbf{Genetic drift} is random change in genetic composition caused by
finite random sampling. Even if the two types have equal reproductive
weights, the number of A individuals can move up or down by chance.

Selection and drift operate in the same run. We do not turn randomness off
when we give A an advantage. Instead, we change the probabilities used by
the random experiment.

\begin{workedexample}
Suppose there are five individuals, and only one is A. A has a reproductive
advantage. Nevertheless, a B individual might be selected as the next parent,
and the one A individual might be selected for replacement. If that happens,
the population becomes all B. The advantage did not disappear. The A lineage
disappeared before reproducing successfully.
\end{workedexample}

\section{Describe the population with a few symbols}

Mathematical symbols save us from writing long sentences repeatedly. Every
symbol below has a concrete meaning. A symbol is not a new biological object.

\subsection{The population size}

The capital letter \(N\) means the total number of individuals. We assume
\(N\geq 2\). The sign \(\geq\) means ``greater than or equal to''. Thus our
population contains at least two individuals.

We keep \(N\) fixed during a simulation. For example, \(N=5\) means the
population always contains five individuals.

\subsection{The number of A individuals}

The lowercase letter \(i\) means the current number of A individuals. It can
take the whole-number values
\[
  i=0,1,2,\ldots,N.
\]
The dots mean that we continue the same counting pattern. The number of B
individuals is \(N-i\): the total minus the number that are A.

For \(N=5\) and \(i=2\), there are two A individuals and three B individuals.
For \(i=0\), A has been lost. For \(i=N\), A has fixed.

The complete state of this model is the count \(i\). We do not need the
positions, ages, names, or family histories of the individuals. Our rules
will depend on the current count and the fixed parameters only.

\subsection{The frequency of A}

The letter \(x\) means the fraction of the population that is A. This is
called the \textbf{allele frequency}. We calculate it by dividing the A count
by the population size:
\[
  x=\frac{i}{N}.
\]
The fraction bar means division. A frequency of \(0\) means no individuals
are A. A frequency of \(1\) means all individuals are A.

For example, two A individuals in a population of five give
\[
  x=\frac{2}{5}=0.4=40\%.
\]
The count and frequency describe the same population in different units.
The count is two individuals. The frequency is forty percent.

\subsection{Reproductive weight and selective advantage}

We give each B individual a reproductive weight of \(1\). The letter \(s\)
means the relative reproductive advantage assigned to A. In our main study,
\(s\geq 0\). Each A individual has reproductive weight \(1+s\).

We also use the letter \(r\) for that same A weight:
\[
  r=1+s.
\]
The two symbols are useful for different purposes. The symbol \(s\) makes
the size of the advantage easy to describe. The symbol \(r\) makes some
algebra shorter.

\begin{workedexample}
If \(s=0.5\), then \(r=1+0.5=1.5\). One A individual has one and a half
times the reproductive weight of one B individual. We can also say that its
weight is fifty percent larger.

This does \emph{not} mean that A has a fifty percent probability of fixation.
It does not even mean that the next parent has a fifty percent chance of
being A. The chance of choosing an A parent also depends on how many A and
B individuals are present.
\end{workedexample}

When \(s=0\), we have \(r=1\). Both types have equal reproductive weight.
We call this the \textbf{neutral} case. Genetic drift still occurs.

\section{Parameters and the changing state}

A \textbf{parameter} is a setting chosen for a run. Here \(N\) and \(s\)
are parameters. Once the run starts, their values stay fixed. The value
\(r=1+s\) is fixed too.

The \textbf{state} is the current count \(i\). That count changes as the run
proceeds. The frequency \(x=i/N\) changes with it.

We must also choose an initial count. For the main experiment, we start with
one A individual. This is a choice about the starting state, not a claim
that one A individual remains throughout the run.

\begin{center}
\begin{tabular}{lll}
\hline
Symbol & Meaning & Example \\
\hline
\(N\) & Total number of individuals & \(5\) \\
\(i\) & Current number of A individuals & \(2\) \\
\(N-i\) & Current number of B individuals & \(3\) \\
\(x\) & Current A frequency, \(i/N\) & \(0.4\) \\
\(s\) & Relative reproductive advantage & \(0.5\) \\
\(r\) & A reproductive weight, \(1+s\) & \(1.5\) \\
\hline
\end{tabular}
\end{center}

\section{The assumptions we agree to use}

Our population is \textbf{well mixed}. Every individual can be selected for
replacement. We do not assign individuals to different locations or groups.
There is no migration into or out of the population.

At each model event, one parent produces one offspring. That offspring has
the same type as its parent. One individual is selected for replacement,
and the offspring takes that individual's place. There is no net change in
population size. The selected individual can be the parent itself. We will
explain why that is allowed when we derive the event probabilities.

The reproductive weights are constant. A does not become more or less
advantageous as the environment changes, because this model has no changing
environment. Its advantage also does not depend on its frequency.

These assumptions allow an exact answer. They also limit what the answer
can tell us about a real population. A careful report states both points.

\section{What one event means for time}

One \textbf{event} is one attempt at reproduction and replacement. Some events
increase the A count. Some decrease it. Some leave it unchanged.

We count unchanged events too. They are still reproduction and replacement
attempts. Omitting them would change the meaning of elapsed time.

For plotting, we may divide the number of events by \(N\). We call this time
in \emph{generations}. It is a conventional time unit: one generation means
\(N\) event attempts. It does not mean that each individual reproduced
exactly once, or that all individuals were replaced simultaneously.

\begin{checkpoint}
\textbf{Questions.} A run has \(N=10\), starts with \(i=3\), and uses
\(s=0.2\). How many B individuals are present? What is the A frequency? What
is the A reproductive weight? What do the two absorbing states mean? If
twenty events have occurred, how many model generations have passed?

\textbf{Worked answers.} There are \(N-i=10-3=7\) B individuals. The A
frequency is \(x=i/N=3/10=0.3\), or thirty percent. The A weight is
\(r=1+s=1+0.2=1.2\). At \(i=0\), A is lost and all ten individuals are B.
At \(i=10\), A is fixed and all ten are A. Twenty events divided by ten
individuals gives two model generations.
\end{checkpoint}

\section{How to use this guide}

Read one worked example before moving to its general equation. Write down
the numbers yourself. When a new symbol appears, say its meaning aloud.
For example, read \(i/N\) as ``the number of A individuals divided by the
total number of individuals''.

When we reach the code, distinguish three questions: what biological event
does this instruction represent, what mathematical quantity does it compute,
and what does Julia do when it executes the instruction? Understanding all
three connections is more useful than memorising the code.

You do not need to understand every later chapter immediately. At this point,
the essential idea is simple: we will follow a count, change it using explicit
random rules, and ask which boundary it eventually reaches.
